\section{Duonic pressure, relation motion, and inertial burden}\label{sec:duonic-pressure-inertia}

Duonic pressure is scalar retained burden. Lead/lag is relation motion. The field attention vector combines them as pressure-weighted lead/lag relation motion.

\begin{definition}[Channel pressure]
For a returned duon-current channel \(e\), define the channel pressure
\[
p^D_e(\partial F,\omega)
=
C_{3,e}
\min\!\left(D_e(\partial F,\omega),|\omega|_e\right).
\]
Here \(|\omega|_e\) is the overlap of the declared window with channel \(e\), and \(D_e\) is the PADAR duration coordinate.
\end{definition}

\begin{definition}[Relation-motion contribution]
For the same channel, define the signed lead/lag relation-motion contribution
\[
r^{(q)}_e(\omega)=a_e(\omega)\Delta^{\chi\mathrm{lag}}_e(\omega)s_{e,q}.
\]
Here \(a_e\) is exposure/retention, \(\Delta^{\chi\mathrm{lag}}_e\) is chiral lead--lag mismatch, and \(s_{e,q}\in\{-1,0,+1\}\) is signed incidence.
\end{definition}

\begin{definition}[Field attention vector]
The field attention vector is pressure-weighted relational motion. In vector form,
\[
\mathbf A_{\mathcal F}(\partial F,\omega)=
\sum_{e\in\mathcal F(\partial F,\omega)}
C_{3,e}\min\!\left(D_e(\partial F,\omega),|\omega|_e\right)
\Delta^{\chi\mathrm{lag}}_e(\omega)\mathbf i_e,
\]
where \(\mathbf i_e\) carries boundary incidence and orientation. Componentwise,
\[
A_{\mathcal F}^{(q)}(\partial F,\omega)=
\sum_{e\in\mathcal F(\partial F,\omega)}
p^D_e(\partial F,\omega)r^{(q)}_e(\omega).
\]
Equivalently,
\[
A_{\mathcal F}^{(q)}(\partial F,\omega)=
\sum_e
C_{3,e}
\min\!\left(D_e(\partial F,\omega),|\omega|_e\right)
 a_e(\omega)\Delta^{\chi\mathrm{lag}}_e(\omega)s_{e,q}.
\]
\end{definition}

The field attention vector records relational motion of the field. Its direction is carried by boundary orientation, incidence, chirality, and signed lead--lag bias. Its magnitude is carried by duonic pressure, retained duration, and PADAR burden. A projection of the field attention vector is a component.

\begin{definition}[Duonic pressure]
The scalar pressure of a field is
\[
P^D_{\mathcal F}(\partial F,\omega)=
\sum_{e\in\mathcal F(\partial F,\omega)}p^D_e(\partial F,\omega).
\]
It records how much retained or pending duon/trit burden is carried on the declared boundary/window.
\end{definition}

When \(P^D_{\mathcal F}(\partial F,\omega)>0\), the normalized attention orientation is
\[
\theta_{\mathcal F}^{(q)}(\partial F,\omega)=
\frac{A_{\mathcal F}^{(q)}(\partial F,\omega)}{P^D_{\mathcal F}(\partial F,\omega)}.
\]
Thus \(P^D\) records scalar burden, \(\theta\) records orientation, and \(A\) records pressure-weighted relation motion.

\paragraph{Inertial burden.}
Duonic pressure resists change of attention orientation. A perturbing duon current changes normalized orientation in inverse proportion to retained pressure:
\[
\Delta\theta_{\mathcal F}^{(q)}=
\frac{\Delta J_{\mathrm{duon}}^{(q)}}{P^D_{\mathcal F}}.
\]
The full field attention vector is then
\[
A_{\mathcal F}^{(q)}=P^D_{\mathcal F}\theta_{\mathcal F}^{(q)}.
\]

\paragraph{Fly / elephant split.}
Two fields can share the same normalized attention orientation while carrying different duonic pressure. The higher-pressure field carries a larger field attention vector because the same orientation is backed by more retained duon burden.
